
\begin{proposition}[A contracted horizontal section is elliptic]\label{cdp:st:C.2.CR.5}
In the situation of \ref{cdp:st:C.2.CR.1} one has
\[
                              g(B)=1.              \tag{C.CR.19}\label{cdp:eq:C.CR.19}
\]
In particular the remaining exact Euler expression is
\[
                            e(S)=n-v+N_0.          \tag{C.CR.20}\label{cdp:eq:C.CR.20}
\]
\end{proposition}
\begin{proof}
By \ref{cdp:st:C.2.CR.2} the only other possibility is $g=2$.
The conductor image is then smooth and reduced at
$q=h(C)$, and its generic transversal is a node or
a cusp. \ref{cdp:st:C.2.CR.3} excludes the node. For the cusp,
the reduced critical locus near $q$ is precisely
that smooth conductor curve: every positive-dimensional
singular locus of a reduced hypersurface surface
belongs to its nonnormal locus, and there are no
other isolated points after the neighbourhood is
shrunk. \ref{cdp:st:C.2.CR.4} makes the germ a product cusp.
Its normalization is smooth, contrary to the
exceptional genus-two curve contracted to the
normalization point $p$. Thus $g=2$ is excluded.
Substitution in (C.CR.3) gives \eqref{cdp:eq:C.CR.20}.
\end{proof}


